% Prospective assistance/retention protocol; no biological experiment reported.
\section{Catch it again when the help is gone}

A computer could guide a hand to catch a ball. The catch happened. What would
happen on the next attempt, after the teaching help stopped? That is the
Brain2Body question. Keep the task and the necessary access to movement
comparable, or the second attempt cannot tell us what the trainee retained.

The instruments introduced in \emph{From a reach to a testable model} would
have to agree on the catch. Scalp EEG records scalp electrical activity,
surface-EMG sleeves on both arms record muscle electrical activity, and the
synchronized, calibrated left/right tripod views would measure the hands and
arms. OpenBCI or an equivalent configurable system remains a candidate whose
sites, capacity, units, delays and clocks need qualification. The question now
is whether those observations can distinguish a successful catch from the
movement that merely looked close.

These left/right hand views differ from the earlier proposed top and bottom
dot maps. Both belong to the motion-capture relay-composite idea, but each
configuration needs its own calibration, timing and visibility checks. A point
the camera cannot see must not quietly become a measured point because a model
can draw it. The virtual body also needs appropriate anatomy and mechanics.
EMG is neither a force measurement nor a recording of incoming sensory nerves;
contact and force claims need observations that can check them.

\subsection{Let each action add a problem worth solving}

Begin by rubbing the hands rapidly together. Touch and movement happen at once.
Matched conditions would help tell their effects apart. The requested knuckle
movement and the stretch with interlaced fingers, palms outward and arms
extended would be recorded as distinct actions. The cold-hands example makes
rubbing familiar; it establishes no cold-exposure result or measured sensation.

Next, pass a juggling ball between the hands in the proposed hot-potato task,
then progress to the preferred two-ball juggling task. Release makes a new
problem for the catching hand. Check the model against the independently
observed handover and its timing, including what remains uncertain. A prediction
from the cue or previous movement might explain success without useful neural
information, so those comparisons stay in the test.

Boxing would ask the system to distinguish stated punch types from measured
movement. NeuroBoxer is the proposed virtual game built around that comparison.
Hands and wrists help stabilize punches as well as the arms moving. A prime-era
Mike Tyson opponent is a proposed virtual game character. The game supplies no
physical opponent or evidence about safe physical boxing. Expert footage can show
movement and tactics within its rights and camera limits. It cannot provide
that expert's brain activity, muscle activity or force measurements.

Toy archery would add another form of coordination: one hand holds while the
other draws, then the fingers release and the hands return. The proposal uses
arrows with small, cup-shaped suction tips and an over-arm pouch with several
arrows. NeuroArcher would use the coordination
against a virtual target. Hitting it cannot, by itself, establish what happened
in a biological nervous system.

The complete snap described in \emph{From a reach to a testable model} tests another
relationship: preparation, contact and return must all agree with the specified
movement. The cup reach remains the simpler precision comparison. Both need
organism-specific anatomy, contact, timing and real-reference evidence; a
human-looking animation cannot establish that a monkey can perform the snap.

\subsection{A virtual hand is the first control milestone}

The first game milestone would follow measured body motion. A later one would
ask whether permitted EEG observations can control a stated virtual action
without using withheld muscle or movement targets. Looking hands-free cannot
settle it. Eye activity, small covert muscle contractions, residual movement,
shared cues and future information need declared controls. Keep the original
brain-only, body-only and paired comparisons, including the simple prediction
that a target stays as it was, and keep the original success rules.

Understanding the intermediary should let us connect a stated intention with
an action as closely as possible for a defined task. The test needs a stated
person, condition, allowed context, tolerated error and independent observation.
It does not establish a universal one-to-one brain map. Predicting a recorded
EEG feature from body observations is likewise a conditional estimate. It does
not identify a unique mind or deliver a sensation.

\subsection{Teaching help and a movement bridge serve different roles}

Remove trainer information to test what was retained. Keep any movement bridge
needed to perform the task, and keep the comparison conditions meaningful.
Measure performance now, retention later and transfer to specified unfamiliar
conditions separately. Report how much help remains and what else could explain
the result. Removing vision in a proposed blindfolded condition changes feedback;
it does not make a test independent or prove learning by itself.

That retention comparison also depends on how movement or feedback reaches the
trainee. EEG and surface EMG supply recordings. A future body writer would need
a separate qualified interface: neuromuscular electrical stimulation (NMES)
and functional electrical stimulation (FES) can cause contractions or functional
movement, while a sensory writer or transducer has a different delivery role.
What the interface delivers would need its own measurement. Physical spinal
work remains read-only. These physical additions are still unqualified and
outside the acquisition and virtual-work budget scope; no such participant
recording, stimulation or actuation is reported here.

Training the computer model and teaching a biological trainee are different
questions. Qualified recordings and simpler baselines come before paired fitting
and fine-tuning. Reinforcement learning would need specified actions, resulting
transitions and rewards; a good reward inside a simulation is not a biological
teaching result. Final held-out observations stay outside training and selection,
and a frozen running model stays separate from a candidate adaptive update. The
current LINK comparison failed both model-admission gates and all five benefit
claims. This task sequence supplies future comparisons, not a new result about
retained skill, repaired tissue or a cure.
